\chapter{Genre and Orchestration: Eight Real-World Ways to Write}
\label{ch:genres}

\classquote{易穷则变，变则通，通则久。}{"Book of Changes", "Appended Phrases II"}{When a way is exhausted, it changes; change makes it flow; flowing, it lasts. Genres stay alive by borrowing and mutating, not by staying pure. A genre is not a label stuck onto a sound. It is a set of choices: which instruments, how the voices divide the work, and where the accents fall. Learning one set of choices is more useful than memorising ten presets.}

\section{Why genres get their own chapter}
\label{sec:genres-why}

The first nine chapters split the subject by function and by instrument family: rhythm, harmony, melody, strings, woodwinds and brass, electronics, and vocals. That is the right way to build concepts, but real songs make all of these choices at once. The rhythm section of a rock song determines how the guitar is written; the sub layer of an EDM track determines how the bass and kick make room for each other.

This volume makes genre content Chapter 11 for three reasons:

\begin{enumerate}
\item \textbf{No repetition}: spreading eight genres across the first ten chapters would repeat the same ``instrumentation--groove--harmony--form'' discussion in different places, and the reader could not compare them side by side.
\item \textbf{Lookup}: when writing rock, you do not need to read the jazz and EDM sections first. Each section in this chapter can stand alone as a recipe.
\item \textbf{Not a substitute for basics}: a genre is a set of choices, not a rule. Read Chapters 1--5 first for the four layer roles, form, groove, harmony and melody; this chapter puts those choices into concrete stylistic skeletons.
\end{enumerate}

\section{Four rulers of orchestration}
\label{sec:orchestration-rulers}

Before the individual genres, sharpen four tools. They apply to every style and give players a shared language for ``why this instrumentation''.

\subsection{Range and transposition}
\label{sec:range-transposition}

\term{range}{range} is the lowest to highest note an instrument or voice can produce; \term{tessitura}{tessitura} is the part of that range where it is most comfortable and usually lives. When orchestrating, ask first ``where should this role sit'', then ``can this instrument reach there''.

\begin{longtable}{p{0.22\textwidth}p{0.30\textwidth}p{0.38\textwidth}}
\toprule
Instrument/voice & Common written range & Transposition and notes \\
\midrule
\endfirsthead
\toprule
Instrument/voice & Common written range & Transposition and notes \\
\midrule
\endhead
Violin & G3--E7 & Non-transposing; bright high register, weak as a long low pad \\
Viola & C3--A6 & Non-transposing; thick middle, good for inner parts \\
Cello & C2--C6 & Non-transposing; can anchor the bass or sing a melody \\
Double bass & E1--G4 & Sounds an octave lower; often written at sounding pitch \\
Flute & C4--C7 & Non-transposing; piercing high, breathy low \\
Oboe & B-flat3--F6 & Non-transposing; focused, good for solo lines \\
Clarinet (B-flat) & Written E3--C7 & Sounds a major second lower; check whether the sample library uses written or concert pitch \\
Bassoon & B-flat1--F5 & Non-transposing; low woodwind, avoid frequent high leaps \\
Trumpet (B-flat) & Written F-sharp3--C6 & Sounds a major second lower; save the high register \\
Horn (F) & Written C3--C6 & Sounds a perfect fourth lower; usually written in the middle-low register \\
Trombone & E2--B-flat4 & Non-transposing; the backbone of low brass \\
Voices & See below & Cast by role, not by extremes \\
\bottomrule
\end{longtable}

Voice starting points: soprano C4--A5, mezzo G3--E5, tenor C3--A4, bass E2--E4. In a real song, make the singer's two or three most comfortable notes ``home''; in the chorus, borrow a third or a fifth above, but do not push the whole song to the edge.

\term{Transposition}{transposition} has two meanings: notated transposition (B-flat clarinet, B-flat trumpet, F horn) and moving a whole line for register or colour. Yeban MIDI sounds at concert pitch by default. If you load a sample library built at written pitch, convert according to its documentation; do not type the written pitch straight into the piano roll. The safe habit is to hear the actual pitch in the piano roll first, then decide whether to transpose.

\subsection{Voice density and space}
\label{sec:density-space}

\term{Voice density}{voice-density} is how many independent parts are saying different things at the same moment; \term{space}{space} is the deliberate gap left unwritten. They are two sides of one coin: density is only audible because of space, and space only has contrast because of density.

Three rules of thumb:

\begin{itemize}
\item Keep one or two ``speaking'' parts per function layer at a time; the rest hold long notes, repeat, or drop out.
\item Low, middle and high registers each need a principal voice; do not put the accent of all three layers on the same beat.
\item A chorus is not ``every track enters''. It is ``every layer is more complete than in the verse''. Save one track for the bridge or outro; it will create more space than filling everything at once.
\end{itemize}

\subsection{Counterpoint and call-and-response}
\label{sec:counterpoint-call}

\term{Counterpoint}{counterpoint} is two or more melodic lines moving at once, each independent yet coordinated. \term{Call-and-response}{call-and-response} is one phrase asking and another answering. Neither belongs only to classical music: a pop counter-melody, a jazz exchange, and an EDM lead/bass dialogue all use the same listening logic.

Do three things before writing counterpoint:

\begin{enumerate}
\item \textbf{Separate registers}: put the lead in the middle-high range and the counter-line lower; do not let both fight in the same band.
\item \textbf{Separate rhythms}: if the lead sings long phrases, answer with short points or offbeats; do not land both on every strong beat.
\item \textbf{Separate directions}: when one line rises, prefer the other falling or holding; parallel same-direction motion easily smears into one line.
\end{enumerate}

Call-and-response is simpler: the call leaves a ``question mark'' (a short rest or an unresolved note), and the answer catches it in the same register or an octave lower. The answer need not be long; two or three notes are often enough.

\subsection{Texture types}
\label{sec:texture-types}

\term{Texture}{texture} is how the parts combine in register, rhythmic density and timbre. Four common types:

\begin{longtable}{p{0.18\textwidth}p{0.34\textwidth}p{0.38\textwidth}}
\toprule
Texture & What it sounds like & Typical use \\
\midrule
\endfirsthead
\toprule
Texture & What it sounds like & Typical use \\
\midrule
\endhead
Homophony & One lead melody plus chordal accompaniment; clear hierarchy & Pop, rock and folk verses \\
Polyphony & Two or more melodies moving independently & Jazz exchanges, string counterpoint, film themes in dialogue \\
Stratified texture & Low, body, air and rhythm each occupy a layer & EDM, film score, pop choruses \\
Ostinato & A short figure repeats; variation comes from displacement or ornament & Rock riffs, electronic ostinati, guofeng loops \\
\bottomrule
\end{longtable}

\section{Eight genres}
\label{sec:genres-eight}

Every section follows the same order: typical instrumentation, rhythm and groove, harmony and melody, form, signature devices, and a Yeban skeleton. Track counts are teaching starting points, not laws; the function division is what must hold.


\section{Side-by-side comparison and exercise}
\label{sec:genres-compare}

Put the eight genres in one table and the differences become clearer:

\begin{longtable}{p{0.13\textwidth}p{0.25\textwidth}p{0.25\textwidth}p{0.25\textwidth}}
\toprule
Genre & Time skeleton & Harmony/melody focus & Most common mistake \\
\midrule
\endfirsthead
\toprule
Genre & Time skeleton & Harmony/melody focus & Most common mistake \\
\midrule
\endhead
Pop & Snare 2/4, kick 1/3 & Clear hook, space at phrase ends & Every track full in the chorus \\
Rock & Clear backbeat, riff-driven & Power chords, short repetition & Guitar and bass on identical notes \\
EDM & Four-on-the-floor, eight-bar blocks & Loop motive, low-end impact & Kick/sub/bass all fighting for lows \\
Hip-hop & Boom-bap/trap pocket & Sample loop, short hook & 808 and kick both full \\
Jazz & Swing, walking bass & ii--V--I, comping & Swing written as even eighths \\
Folk & Fingerpicking figure, sparse & I--IV--V--vi, narrative & Full kit destroys the breath \\
Film score & Picture-driven, ostinato & Theme motive, modal colour & Full orchestra all the way, no space \\
Guofeng & Medium, expansive, breaths & Pentatonic, ornament & Pentatonic label, stacked folk instruments \\
\bottomrule
\end{longtable}

\begin{exercise}
Choose a genre you do not know well and build eight bars from its ``Yeban skeleton'':
\begin{enumerate}
\item Create and name tracks, and write section markers first; do not hunt for a preset.
\item Write only the rhythm section and harmony layer; solo the group and confirm the low end is not fighting.
\item Add the melody/vocal and one colour layer; mute each new layer once and confirm it actually helps.
\item Finish by answering: which track, if removed, would make the music worse? If you cannot answer, that track may be only a habit.
\end{enumerate}
\end{exercise}

\section{Chapter takeaways}
\label{sec:genres-takeaways}

\begin{itemize}
\item A genre is a combination of instrumentation, groove, harmony, form and signature devices, not a timbre label.
\item Range, transposition, density, space, counterpoint and texture are the four rulers shared by all eight styles.
\item In every style, protect the function division before discussing instrument count; few and clear usually beats many and muddy.
\item A Yeban skeleton lands in the order ``create tracks--write skeleton--add layers--solo check''; do not start by stacking presets.
\item Styles may borrow devices, but must not borrow another style's signature. Before borrowing, ask: does it still recognise itself?
\end{itemize}

\section{Further reading}
\label{sec:11-genres-further-reading}
\begin{itemize}
  \item \href{https://en.wikipedia.org/wiki/Music_genre}{Music genre}: See how genres are classified by instruments, texture and use.
  \item \href{https://en.wikipedia.org/wiki/Arrangement}{Arrangement}: Return to the choices rather than memorising presets.
  \item \href{https://www.youtube.com/watch?v=M6WhsnWIhcM}{Defining a style through layers and profiles}: How layers and instrument profiles combine into a style.
  \item \href{https://www.youtube.com/watch?v=uO4_OuZHh1Q}{One kit, different genre grooves}: Compare how drum writing changes across genres.
\end{itemize}

\begin{helpmapping}
\begin{longtable}{p{0.45\textwidth}p{0.45\textwidth}}
\toprule
Section & Help anchor \\
\midrule
Why genres get their own chapter & \texttt{arrangement-genres-why} \\
Range and transposition & \texttt{arrangement-genres-range} \\
Voice density and space & \texttt{arrangement-genres-density} \\
Counterpoint and call-and-response & \texttt{arrangement-genres-counterpoint} \\
Texture types & \texttt{arrangement-genres-texture} \\
Side-by-side comparison and exercise & \texttt{arrangement-genres-compare} \\
\bottomrule
\end{longtable}
\end{helpmapping}
